\documentclass[11pt,a4paper]{article}

% --- Packages ---
\usepackage{amsmath}
\usepackage{amssymb}
\usepackage{geometry}
\usepackage{graphicx}
\usepackage{hyperref}
\usepackage{siunitx}

% Page layout
\geometry{margin=1in}

% Hyperlink colors
\hypersetup{
    colorlinks=true,
    linkcolor=blue,
    urlcolor=blue
}

\title{Reinforcement Learning for Sun Tracking}
\author{}
\date{}

\begin{document}

\maketitle

\begin{abstract}
This document describes the simulation environment used to train a reinforcement learning (RL) agent for solar tracking. 
It explains how the Sun is modeled, how the dish moves, and how the power signal is computed. 
The goal is to create a simple but realistic setup in which an RL algorithm can learn to point a dish toward the Sun.
\end{abstract}

\section{Introduction}

This project explores a simple reinforcement learning (RL) approach for solar tracking.  
The goal is for an agent to learn how to orient a dish so that it points towards the center of the Sun, using only power sensor readings and its own motion commands.  
We build a simulation that includes:
\begin{itemize}
    \item a model of the Sun's position over time,
    \item a basic physical model of the dish,
    \item a power measurement based on how well the dish aligns with the Sun,
    \item a Gym-style environment for training an RL agent.
\end{itemize}

The reward is the measured power coming from the Sun, so the agent learns to move the dish in a way that maximizes this value.  
The following sections describe the simulation components that make up the environment.


\section{Sun Model}

The Sun's position is computed using the Astral library, based on a geographic location and the current UTC time.  
Astral provides the Sun's azimuth (degrees clockwise from North) and elevation (degrees above the horizon).  
These angles are then converted into a unit vector expressed in a local East--North--Up (ENU) coordinate frame.

\subsection{Conversion to ENU Vector}

Given an azimuth angle $\mathrm{az}$ and an elevation angle $\mathrm{el}$ (both in radians), the corresponding unit vector
$\mathbf{s} = (s_x, s_y, s_z)$ is computed as:
\begin{equation}
\label{eq:enu}
s_x = \cos(\mathrm{el}) \cos(\mathrm{az}), \qquad
s_y = \cos(\mathrm{el}) \sin(\mathrm{az}), \qquad
s_z = \sin(\mathrm{el}).
\end{equation}

This vector points from the dish location toward the Sun at the given moment in time(our point of reference in the simulation is the dish).

\subsection{Time Update}

The simulation keeps track of time in seconds.  
At each step, the current UTC time is:
\begin{equation}
t_{\text{now}} = t_{\text{start}} + \Delta t,
\end{equation}
where $\Delta t$ is the elapsed simulation time.  

The Sun vector $\mathbf{s}(t)$ is recomputed at every step using Astral and the
ENU conversion formula shown in Equation~\ref{eq:enu}.

\section{Dish Dynamics Model}

The dish is modeled with two angles: azimuth $\mathrm{az}$ and elevation $\mathrm{el}$.  
Each angle has an angular velocity, $\dot{\mathrm{az}}$ and $\dot{\mathrm{el}}$.  
The reinforcement learning agent controls the dish by outputting two values, one for azimuth angular velocity and one for elevation angular velocity.  
These values are interpreted as the \emph{commanded angular velocities} that the motors should try to reach.

\subsection{Motor Lag Model}

Let $v_{\text{cmd}}$ be the commanded angular velocity (either in azimuth or elevation) produced by the RL agent.  
Because real motors cannot change speed instantly, the actual angular velocity follows a simple first-order response:

\begin{equation}
\label{eq:lag}
x \leftarrow x + \left( v_{\text{cmd}} - x \right) \frac{\Delta t}{\tau},
\end{equation}

where $x$ represents $\dot{\mathrm{az}}$ or $\dot{\mathrm{el}}$,  
$\Delta t$ is the simulation time step,  
and $\tau$ is the motor time constant.

\subsection{Angle Integration}

After the angular velocities are updated, the angles themselves are integrated:

\begin{equation}
\label{eq:az_update}
\mathrm{az} \leftarrow \mathrm{wrap}\left( \mathrm{az} + \dot{\mathrm{az}} \, \Delta t \right),
\end{equation}

\begin{equation}
\label{eq:el_update}
\mathrm{el} \leftarrow \mathrm{clip}\left( \mathrm{el} + \dot{\mathrm{el}} \, \Delta t \right).
\end{equation}

Azimuth is wrapped to stay within $[-\pi, \pi]$,  
while elevation is clipped so the dish does not point below the horizon.

\subsection{Dish Pointing Vector}

The dish orientation is represented as a unit vector $\mathbf{n}$ in the East--North--Up frame:

\begin{equation}
\label{eq:n_vector}
n_x = \cos(\mathrm{el}) \cos(\mathrm{az}), \qquad
n_y = \cos(\mathrm{el}) \sin(\mathrm{az}), \qquad
n_z = \sin(\mathrm{el}).
\end{equation}

This vector defines the direction the dish is currently facing.

\section{Power Measurement Model}

The dish receives power based on how well it is aligned with the Sun.  
Both the dish direction $\mathbf{n}$ and the Sun direction $\mathbf{s}$ are represented as unit vectors in the East--North--Up frame.  
The measured power is computed using their dot product.

\subsection{Ideal Power Without Noise}

The ideal (noise-free) power value is:

\begin{equation}
\label{eq:power_dot}
P = \max\left(0,\; \mathbf{n} \cdot \mathbf{s} \right),
\end{equation}

which ensures that power is zero when the dish points away from the Sun.  
The dot product reaches $1$ when the dish is perfectly aligned with the Sun(dish looking exactly at the sun's center).

\subsection{Power Sensor Noise}

A small amount of noise is added to simulate measurement uncertainty.  
If the noise standard deviation is $\sigma$, then the final measured power is:

\begin{equation}
\label{eq:noise}
P \leftarrow \mathrm{clip}\left( P + \varepsilon \right),
\qquad
\varepsilon \sim \mathcal{N}(0, \sigma^2),
\end{equation}

where the value is clipped to stay within the range $[0, 1]$.  
This noise term helps make the environment slightly more realistic and prevents the agent from relying on perfectly smooth measurements.

\section{Reinforcement Learning Environment}

The simulation is wrapped in a Gym-style environment so that an RL agent can interact with it.  
At each step, the agent observes the current state, outputs its action, and receives a reward based on how well the dish aligns with the Sun.

\subsection{Observation Space}

The environment provides a six-dimensional observation vector:

\begin{equation}
o = 
\left[
\mathrm{az},\;
\mathrm{el},\;
\dot{\mathrm{az}},\;
\dot{\mathrm{el}},\;
P,\;
\Delta P
\right],
\end{equation}

where:
\begin{itemize}
    \item $\mathrm{az}$ and $\mathrm{el}$ are the dish angles,
    \item $\dot{\mathrm{az}}$ and $\dot{\mathrm{el}}$ are the angular velocities,
    \item $P$ is the measured power,
    \item $\Delta P = P - P_{\text{prev}}$ is the change in power.
\end{itemize}

The value $\Delta P$ helps indicate whether the dish is moving toward or away from the Sun.

\subsection{Action Space}

The agent outputs two continuous values:

\begin{equation}
a = [a_{\text{az}},\; a_{\text{el}}], \qquad a_{\text{az}}, a_{\text{el}} \in [-1, 1].
\end{equation}

These values do not represent angles.  
Instead, each one specifies how fast the agent wants the dish to rotate.  
They are converted into commanded angular velocities by scaling with the maximum motor speed $v_{\max}$:

\begin{equation}
v_{\text{cmd}} = a \cdot v_{\max}.
\end{equation}

For example, an action of $1.0$ means ``rotate at the maximum speed'',  
$0.0$ means ``do not rotate'',  
and $-1.0$ means ``rotate at maximum speed in the opposite direction''.  
The commanded velocities are then used in the motor lag model.

\subsection{Reward Function}

The reward at each step is the measured power:

\begin{equation}
r = P.
\end{equation}

Maximizing this reward encourages the agent to orient the dish so that it points directly at the Sun.

\subsection{Episode Termination}

Episodes run for a fixed amount of simulated time.  
Let $T_{\max}$ be the maximum allowed duration.  
After every simulation step, the elapsed time increases by $\Delta t$.  
The episode continues as long as:

\begin{equation}
t_{\text{episode}} \le T_{\max}.
\end{equation}

Once the elapsed time exceeds $T_{\max}$, the episode ends.  
There is no separate success or failure condition; the goal is simply to track the Sun over this fixed time period.

\subsection{Timing and Control Frequency}

Each environment step corresponds to a fixed simulation time step $\Delta t$.  
In this project, we use:
\begin{equation}
\Delta t = 0.5 \text{ seconds}.
\end{equation}

This means the reinforcement learning agent produces a \emph{new action every 0.5 seconds}.  
Each action specifies the desired angular velocities for that 0.5 seconds interval, and the dish dynamics update accordingly.  
The angle change in each step is determined by:

\begin{equation}
\Delta \mathrm{az} = \dot{\mathrm{az}} \cdot \Delta t, \qquad
\Delta \mathrm{el} = \dot{\mathrm{el}} \cdot \Delta t,
\end{equation}

where $\dot{\mathrm{az}}$ and $\dot{\mathrm{el}}$ are the actual motor speeds after applying the motor lag model.  
Thus, the dish motion depends both on the agent's commanded velocities and on the physical time step.

Since episodes last $T_{\max}$ seconds and each step lasts $\Delta t$ seconds, each episode contains:

\begin{equation}
N_{\text{steps}} = \frac{T_{\max}}{\Delta t}
\end{equation}

decision points for the agent.  
For example, with $T_{\max} = 90$ seconds and $\Delta t = 0.5$ seconds, each episode contains 180 steps.

\section{Reinforcement Learning Algorithm}

To learn how to orient the dish toward the Sun, we use the Proximal Policy Optimization (PPO) algorithm from the Stable-Baselines3 library.  
PPO is well suited for continuous control and is stable enough to train on noisy observations.

\subsection{Policy Network}

The policy and value networks are standard multilayer perceptrons.  
They take the six-dimensional observation vector as input and output the two commanded angular velocities (after normalization).

The network architecture is:

\begin{itemize}
    \item Two fully connected hidden layers
    \item 128 neurons per layer
    \item ReLU activation functions
\end{itemize}

This architecture is specified as:
\begin{equation}
\text{net\_arch} = [128, 128].
\end{equation}

\subsection{Training Setup}

The environment is wrapped in a monitor and a vectorized environment wrapper so that PPO can collect rollouts.  
The main hyperparameters are:

The main hyperparameters are:

\begin{itemize}
    \item Learning rate: $3 \times 10^{-4}$
    \item Discount factor: $\gamma = 0.99$
    \item GAE parameter: $\lambda = 0.95$
    \item Rollout length: 1024 steps (this means PPO collects 1024 environment steps before updating the policy)
    \item Batch size: 256 (this is how many samples are used at once during a gradient update)
    \item Entropy coefficient: 0.0
\end{itemize}


The model is trained for:
\begin{equation}
1{,}000{,}000 \text{ timesteps}.
\end{equation}

During training, the sensor noise is enabled \texttt{(noise\_std = 0.01)} to make the policy robust to measurement imperfections.

\subsection{Training Environment Parameters}

Each environment used for training has the following parameters:

\begin{itemize}
    \item Simulation timestep: $\Delta t = 0.5$ seconds
    \item Maximum motor speed: $10^\circ$/s
    \item Motor lag constant: $\tau = 0.3$
    \item Sensor noise: $0.01$
    \item Episode horizon: $T_{\max} = 90$ seconds
\end{itemize}

At each simulation step, the agent outputs a new velocity command for azimuth and elevation.  
The dish motion is integrated over $\Delta t = 0.5$ seconds, so the agent makes 180 decisions per episode.

\subsection{Evaluation Procedure}

After training, the policy is evaluated over multiple episodes using a separate evaluation function.  
For each evaluation episode, we compute:

\begin{itemize}
    \item Total episode reward,
    \item The final power measurement,
    \item The average power over the last $k$ steps (the ``tail average'').
\end{itemize}

The tail average is used because it is more stable than a single noisy final reading.  
A policy is considered successful if the tail average exceeds $0.95$, meaning the dish is almost perfectly aligned during the end of the episode.

\subsection{Multi-Seed Training}

To test the consistency and reliability of the PPO method, we repeat the full training process with several different random seeds.  
For each seed:

\begin{enumerate}
    \item A new PPO model is trained for 1,000,000 steps.
    \item The model is saved to disk.
    \item The model is evaluated in two conditions:
    \begin{itemize}
        \item Clean evaluation (\texttt{noise\_std = 0.0}),
        \item Noisy evaluation (\texttt{noise\_std = 0.01}).
    \end{itemize}
\end{enumerate}

Each evaluation runs over many episodes (typically 200) to obtain stable statistics.  
For every seed, we record:

\begin{itemize}
    \item Average episode return,
    \item Average tail power,
    \item Success rate (fraction of episodes with tail power $> 0.95$).
\end{itemize}

\subsection{Success Rate Comparison Across Seeds}

After all seeds are trained, the success rates are plotted to show how training stability varies with random initialization.  
The plot shows, for each seed:

\begin{itemize}
    \item Success rate in the clean environment,
    \item Success rate in the noisy environment.
\end{itemize}

This multi-seed experiment provides a clearer picture of how reliable the learned controller is and how sensitive training is to randomness in initialization and environment conditions.




\end{document}
